\documentclass[11pt,a4paper]{article}
\usepackage[margin=2.4cm]{geometry}
\usepackage{amsmath,amssymb,amsthm,mathtools}
\usepackage[protrusion=true,expansion=false]{microtype}
\usepackage{booktabs,enumitem}
\usepackage{placeins}
\usepackage{tikz}
\usetikzlibrary{arrows.meta,positioning,calc,fit,backgrounds}
\definecolor{tierspend}{RGB}{200,60,60}
\definecolor{tierfull}{RGB}{210,140,40}
\definecolor{tierin}{RGB}{60,130,180}
\definecolor{tierrecv}{RGB}{70,150,90}
\usepackage[colorlinks=true,linkcolor=blue!50!black,citecolor=blue!50!black,urlcolor=blue!50!black]{hyperref}

\emergencystretch=3em
\hbadness=10000
\hfuzz=2pt

\newtheorem{theorem}{Theorem}[section]
\newtheorem{lemma}[theorem]{Lemma}
\newtheorem{proposition}[theorem]{Proposition}
\newtheorem{corollary}[theorem]{Corollary}
\theoremstyle{definition}
\newtheorem{definition}[theorem]{Definition}
\newtheorem{construction}[theorem]{Construction}
\newtheorem{assumption}[theorem]{Assumption}
\newtheorem{remark}[theorem]{Remark}
\newtheorem{example}[theorem]{Example}

\newcommand{\F}{\mathbb{F}}
\newcommand{\Z}{\mathbb{Z}}
\newcommand{\N}{\mathbb{N}}
\newcommand{\G}{\mathbb{G}}
\newcommand{\E}{\mathbb{E}}
\newcommand{\PP}{\mathbb{P}}
\renewcommand{\Pr}{\mathrm{Pr}}
\newcommand{\Fp}{\F_p}
\newcommand{\Fq}{\F_q}
\newcommand{\ip}[2]{\langle #1,\, #2 \rangle}
\newcommand{\Comm}{\mathsf{Commit}}
\newcommand{\Open}{\mathsf{Open}}
\newcommand{\Setup}{\mathsf{Setup}}
\newcommand{\Verify}{\mathsf{Verify}}
\newcommand{\negl}{\mathsf{negl}}
\newcommand{\poly}{\mathsf{poly}}
\newcommand{\PPT}{\mathsf{PPT}}
\newcommand{\Adv}{\mathcal{A}}
\newcommand{\Prov}{\mathcal{P}}
\newcommand{\Ver}{\mathcal{V}}
\newcommand{\Sim}{\mathcal{S}}
\newcommand{\Ext}{\mathcal{E}}
\newcommand{\Rel}{\mathcal{R}}
\newcommand{\Lang}{\mathcal{L}}
\newcommand{\nfsym}{\mathsf{nf}}
\newcommand{\cmsym}{\mathsf{cm}}
\newcommand{\cmx}{\mathsf{cmx}}
\newcommand{\cvsym}{\mathsf{cv}}
\newcommand{\rk}{\mathsf{rk}}
\newcommand{\Pallas}{\ensuremath{\mathsf{Pallas}}}
\newcommand{\Vesta}{\ensuremath{\mathsf{Vesta}}}
\newcommand{\Ep}{\ensuremath{\mathcal{E}}}
\newcommand{\NoteCommit}{\ensuremath{\mathsf{NoteCommit}}}
\newcommand{\Sinsemilla}{\ensuremath{\mathsf{Sinsemilla}}}
\newcommand{\ShortCommit}{\ensuremath{\mathsf{ShortCommit}}}
\newcommand{\CommitIvk}{\ensuremath{\mathsf{Commit}^{\mathsf{ivk}}}}
\newcommand{\Extract}{\ensuremath{\mathsf{Extract}}}
\newcommand{\Acc}{\ensuremath{\mathsf{Acc}}}
\newcommand{\GroupHash}{\ensuremath{\mathsf{GroupHash}}}
\newcommand{\ask}{\ensuremath{\mathsf{ask}}}
\newcommand{\aksym}{\ensuremath{\mathsf{ak}}}
\newcommand{\nksym}{\ensuremath{\mathsf{nk}}}
\newcommand{\ivk}{\ensuremath{\mathsf{ivk}}}
\newcommand{\fvk}{\ensuremath{\mathsf{fvk}}}
\newcommand{\ovk}{\ensuremath{\mathsf{ovk}}}
\newcommand{\dk}{\ensuremath{\mathsf{dk}}}
\newcommand{\pkd}{\ensuremath{\mathsf{pk_d}}}
\newcommand{\gd}{\ensuremath{\mathsf{g_d}}}
\newcommand{\rcv}{\ensuremath{\mathsf{rcv}}}
\newcommand{\rcm}{\ensuremath{\mathsf{rcm}}}
\newcommand{\rivk}{\ensuremath{\mathsf{r_{ivk}}}}
\newcommand{\rsk}{\ensuremath{\mathsf{rsk}}}

\renewenvironment{abstract}{%
  \small\begin{center}{\bfseries\abstractname}\end{center}\medskip\noindent\ignorespaces
}{\par\medskip}

\title{\textbf{\Huge Voting Guide}\\[6pt]\large The Zally shielded voting protocol, as deployed for Zcash coinholder polling}
\author{m@rek.onl}
\date{}
\begin{document}
\maketitle
\thispagestyle{empty}
\begin{abstract}
This volume of \emph{The Zcash Arboretum} documents the shielded voting protocol
deployed for Zcash coinholder polling. No specification of the protocol has
merged anywhere; the commit-pinned implementation is the de facto
specification, and this volume documents that implementation: ballot structure
and eligibility proving against a snapshot of the note-commitment tree, the
layered nullifier defences against double voting, threshold ElGamal tallying
under a validator DKG, and --- stated without euphemism --- the trust
assumptions and the one confirmed gap between the published description and
the deployed circuit.
\end{abstract}
\clearpage
\tableofcontents
\clearpage
\section{The protocol and its provenance}\label{sec:vote-context}

This volume documents the \emph{Shielded Voting Protocol} --- internally
``Zally'', shipped to users as \emph{Coinholder Polling} --- the first
governance mechanism to run over Zcash's shielded pool in production. The
protocol lets a holder of Orchard notes vote with the weight of their unspent
funds, privately: the funds never move, the notes and addresses are never
revealed, and only aggregate totals become public. The internal name appears
in the client library's own description of itself (``the Zally governance
protocol'', \texttt{valargroup/zcash\_voting} @ \texttt{624cf198},
\texttt{README.md}); the draft ZIP carries the public name
(\texttt{draft-valargroup-shielded-voting.md}, zcash/zips PR \#1200); the
wallet changelog carries the product name (\texttt{zashi-android} @
\texttt{b4edd6f6}, \texttt{CHANGELOG.md}). This section fixes what shipped,
where the design came from, and --- decisive for every citation in this
volume --- what the specification actually is.

\subsection{What shipped}\label{sec:vote-shipped}

In outline: a voting round pins a snapshot of the Zcash chain --- an Orchard
note-commitment root and a nullifier-set root. A holder proves in zero
knowledge that they control notes which existed and were unspent at the
snapshot, and delegates their quantised weight --- one ballot per $0.125$~ZEC
($12{,}500{,}000$ zatoshi per ballot,
\texttt{voting-circuits/src/params.rs}) --- to a round-scoped governance
hotkey. The hotkey casts each vote as ElGamal-encrypted split shares on a
dedicated permissioned chain, \texttt{svoted}, whose ballot surface is five
messages (\texttt{vote-sdk/proto/svote/v1/tx.proto}); relayers reveal shares
with timing decorrelation; and a per-round threshold \emph{Election
Authority} --- a distributed key generation among the chain's validators ---
decrypts aggregates only (\texttt{vote-sdk/x/vote/keeper/keeper\_tally.go}).
The remainder of this volume constructs each of these objects in turn; here
we record only their provenance.

\paragraph{Deployment.} Coinholder Polling shipped in Zodl --- the wallet
formerly Zashi; the team left ECC in January 2026 and the application package
identity is unchanged --- version 3.5.0 on Android and iOS, released
2026-05-28, with signing support for the Keystone hardware wallet
(\texttt{zashi-android} @ \texttt{b4edd6f6}, \texttt{CHANGELOG.md}, entry
\texttt{[3.5.0 (1736)]}); a follow-up release, 3.5.2, followed on 2026-06-04
(ibid., entry \texttt{[3.5.2 (1742)]}). The first production use was the NU7
sentiment coinholder poll, open 2026-06-10 through 2026-06-24 with results
published by 2026-06-29, run in parallel with a ZCAP poll of the same
question (statuses checked 2026-07-08).
%% TODO: pin the poll announcement and results URLs once archived; production
%% coordinator/validator identities not yet established from public data.

\paragraph{Authorship.} Valar Group designed and implemented the protocol,
led by Dev Ojha (GitHub \texttt{ValarDragon}); the draft ZIP is authored by
GitHub user \texttt{p0mvn}, also of Valar Group. The upstream
\texttt{orchard} crate gained the protocol's one in-tree dependency --- the
\texttt{unstable-voting-circuits} feature, which widens circuit internals
(the \NoteCommit{}, nullifier-derivation, and \CommitIvk{} chips) for
out-of-tree reuse --- in a commit by Adam Tucker (zcash/orchard @
\texttt{c0b6b917}, 2026-04-23). ZODL engineers str4d (Jack Grigg),
Daira-Emma Hopwood, and Kris Nuttycombe reviewed --- they did not design ---
the architecture in Q1~2026 (ZODL retroactive grant application,
ZcashCoinholderGrantsProgram issue \#33); ZODL also built the wallet
integration.

\subsection{Lineage}\label{sec:vote-lineage}

Coin-weighted shielded polling on Zcash has one prior implementation:
hhanh00's \texttt{zcash-vote}, a standalone application that ran the 2024
developer-fund poll and the April--May 2025 lockbox poll
(\texttt{hhanh00/zcash-vote} @ \texttt{f05f5f25},
\texttt{src/election.rs}; it depends on a forked \texttt{orchard} at rev
\texttt{75448e67} with a \texttt{vote} feature). Its ZIP --- zcash/zips PR
\#853, ``Coin Voting ZIP'' --- has been open since 2024-05-28 and remains
unmerged as of 2026-07-08. Zally is independent of it: neither codebase
references the other (verified by search over both trees, 2026-07-08), and
Zally reuses the upstream \texttt{orchard} crate through a feature flag
where \texttt{zcash-vote} patches in its fork. The two share the
core idea --- pin a snapshot, prove note eligibility against it, and prevent
double voting with a per-election nullifier --- but Zally adds hotkey
delegation through a dedicated note type (the Vote Authority Note), split
ElGamal shares, a threshold Election Authority, relayers, private
eligibility queries by single-server PIR, a dedicated chain, and hardware
signing, and it lives inside the production wallet rather than beside it.

Two adjacent mechanisms are not ancestors. ZCAP, the Zcash Community
Advisory Panel, polls an identified panel through Helios with no coin
weighting; the June 2026 NU7 question ran through both instruments in
parallel, which makes the distinction operational, not academic. Hopwood's
draft \texttt{draft-ecc-onchain-accountable-voting.md} (created 2025-02-21,
in the merged zips tree at \texttt{c6b70358}) targets the complement ---
accountable on-chain voting by identified key-holders --- and shares no
design lineage with either.

\subsection{The source situation: code as the de facto specification}
\label{sec:vote-source-situation}

No normative document exists. The would-be specification, zcash/zips PR
\#1200 (``[ZIP TBD] Shielded Voting Protocol'', opened 2026-03-05), is open
and unmerged as of 2026-07-08, not updated since 2026-06-02, and carries a
review comment flagging a section as ``Underspecified''; the companion PIR
draft, PR \#1198, is likewise open. The design
book the implementation itself cites --- \texttt{valargroup/shielded-vote-book},
linked from the \texttt{zcash\_voting} README --- is private, as is its
mirror \texttt{z-cale/shielded-vote-book} (both return 404 to anonymous
access, verified 2026-07-08). The public documentation
(\texttt{valargroup.gitbook.io/shielded-vote-docs}, fetched 2026-07-08) is
prose: its circuit pages list public and private inputs and some Poseidon
preimages, but omit the governance-nullifier, VAN-nullifier, and
share-nullifier formulas, the nullifier-domain derivation, the
proposal-bitmask width, and the indexed-Merkle-tree leaf model --- and it
misstates the ballot-quantisation rule, the sole open problem this volume
records (the closing paragraph of this section).

The consequence governs this volume: the commit-pinned implementation of
Table~\ref{tab:vote-sources} is the de facto specification, and what we
document is that implementation. Every governance-specific formula in the
chapters that follow is read from these sources at these commits, never from
designer prose; where the prose and the code disagree, the code is the
protocol, and we flag the disagreement.

\begin{table}[htbp]
\centering
\small
\begin{tabular}{@{}llll@{}}
\toprule
Repository & Commit & Date & Role \\
\midrule
\texttt{valargroup/voting-circuits} & \texttt{4c39abd5} & 2026-06-03 &
  circuits: prover \emph{and} verifier \\
\texttt{valargroup/vote-sdk} & \texttt{cb915f51} & 2026-06-05 &
  vote chain: validation, DKG, tally \\
\texttt{valargroup/zcash\_voting} & \texttt{624cf198} & 2026-06-07 &
  client: proofs, scheduling, persistence \\
\texttt{valargroup/vote-nullifier-pir} & \texttt{603e4e3b} & 2026-06-03 &
  PIR: nullifier non-membership \\
\texttt{valargroup/spendability-pir} & \texttt{a924008d} & 2026-06-04 &
  PIR: spendability and witness queries \\
\texttt{zcash/orchard} & \texttt{c0b6b917} & 2026-04-23 &
  \texttt{unstable-voting-circuits} feature \\
\bottomrule
\end{tabular}
\caption{The de facto specification: the sources this volume documents,
pinned by commit. Wallet-side integration is pinned separately:
\texttt{zcash-android-wallet-sdk} @ \texttt{f386369e} depends on
\texttt{zcash\_voting~=0.11.0} and \texttt{orchard~=0.14.0} with the
\texttt{unstable-voting-circuits} feature
(\texttt{backend-lib/Cargo.toml}); \texttt{zashi-android} @
\texttt{b4edd6f6}, \texttt{zashi-ios} @ \texttt{8ab0e6ed}.}
\label{tab:vote-sources}
\end{table}

\paragraph{The single-crate caveat.} Crate \texttt{voting-circuits} exposes
both the provers and the verifiers of all three circuits
(\texttt{verify\_delegation\_proof}, \texttt{verify\_vote\_proof}, and
\texttt{verify\_share\_reveal\_proof}, each in its circuit directory's
\texttt{prove.rs}), and the
chain consumes the same crate for validation:
\texttt{vote-sdk/circuits/Cargo.toml} pins \texttt{voting-circuits} at
version \texttt{0.8.0}, and \texttt{vote-sdk/x/vote/ante/validate.go}
invokes it. Wallet and chain therefore
cannot drift apart --- but their agreement is a tautology, not evidence of
correctness. No independent second implementation exists that could catch a
mis-transcribed relation, and no merged specification exists against which
either side could be audited. The reader should weigh every constraint we
state in later sections with this in mind.

\paragraph{Classification.} Design-stage material in this volume obeys a
three-bin discipline. Every claim is exactly one of: \emph{specified} --- a
merged normative artifact exists and we cite it; \emph{designed but
unspecified} --- the design is concrete and open-source but no normative
artifact pins it, so we cite the implementation and state what a
specification would need to fix; or an \emph{open problem} --- the design
itself leaves the question unresolved. The bin is always visible in the
prose.

\paragraph{Specified.} Only the inherited primitives: \NoteCommit{},
nullifier derivation, \CommitIvk{}, the \Sinsemilla{} Merkle instance,
RedPallas, the ZIP-244 sighash, and ZIP-32 key derivation --- each specified
in the Zcash protocol specification and the ZIPs, constructed in the
\emph{Ironwood Guide} (\S``Representation of a note: notes and note
commitments'', \S``Prevention of double-spending: nullifiers'', \S``Keys:
capabilities for spending and for viewing'', \S``Proof of ownership of
\emph{some} valid note: the commitment tree'', \S``Authorisation of the
spend: RedPallas''), and reused unchanged through the
\texttt{unstable-voting-circuits} feature (zcash/orchard @
\texttt{c0b6b917}). No governance-specific object has a specification of
that rigor.

\paragraph{Designed but unspecified.} Everything else this volume presents:
the five-message ballot structure, the delegation relation with its
indexed-Merkle-tree non-membership argument, the coordinator-posted snapshot
roots, the Vote Authority Note and its synthetic-action authorisation, all
four double-vote layers, the split-share ElGamal encoding and vote
commitment, the relayer protocol and the trust it grants, the per-round DKG
and chain-verified tally, and the PIR subsystem. For each we cite the pinned
commit and file, and we note what PR \#1200 would need to pin down for the
object to change bins.

\paragraph{Open problems.} One, and only one, was found: the
ballot-quantisation relation enforced by the delegation circuit is strictly
weaker than the exact floor division the public documentation states --- the
relation admits a one-ballot \emph{under}-claim for roughly a third of
note-value totals, while an over-claim remains impossible
(\texttt{voting-circuits/src/delegation/circuit.rs},
\texttt{src/delegation/README.md} \S8). We state and analyse the gap where
delegation weight is constructed. It is the only spec-versus-code
disagreement this volume records.

\section{Ballots and eligibility}\label{sec:vote-ballot}

A voting round is a conversation in five messages on the vote chain, and
admission to it is a zero-knowledge proof against a pinned Zcash snapshot.
This section constructs both sides of that admission: the wire surface of a
round (\S\ref{sec:vote-messages}); the eligibility relation the delegation
circuit proves --- ownership of Orchard notes that existed at the snapshot
and were unspent at it (\S\ref{sec:vote-zkp1}, \S\ref{sec:vote-imt}); the
provenance of the two snapshot roots those proofs are checked against
(\S\ref{sec:vote-snapshot}); and the quantization that turns note value into
voting weight (\S\ref{sec:vote-weight}). Except where marked otherwise,
every claim below sits in the \emph{designed but unspecified} bin of
\S\ref{sec:vote-source-situation}: we read it from the commit-pinned sources
of Table~\ref{tab:vote-sources}, and each subsection closes by recording its
bin. The quantization carries this volume's single open problem; we state
and prove it where it arises.

\subsection{The five-message round}\label{sec:vote-messages}

The ballot surface of the vote chain is five protobuf messages
(\texttt{vote-sdk/proto/svote/v1/\allowbreak tx.proto}). In lifecycle order: a
coordinator opens the round, holders delegate weight, governance hotkeys
cast votes, helper servers reveal shares, and the round closes with a tally.
The three voter-facing messages (2--4) are not ordinary Cosmos transactions:
they bypass the Cosmos transaction envelope entirely, and the chain
authenticates each by its zero-knowledge proof --- messages~2 and~3
additionally by a RedPallas signature; message~4 carries none --- rather
than by an account signature (\texttt{tx.proto:10-13};
\texttt{vote-sdk/x/vote/ante/\allowbreak validate.go}).

\paragraph{Message 1: \texttt{MsgCreateVotingSession}.} The round-opening
message carries the snapshot pin and the agenda: \texttt{snapshot\_height},
\texttt{snapshot\_blockhash}, \texttt{proposals\_hash},
\texttt{vote\_end\_time}, the two ledger anchors
\texttt{nullifier\_imt\_root} and \texttt{nc\_root}, the proposal list, and
display metadata (\texttt{tx.proto:36-48}). It is not a public RPC: it is
absent from the \texttt{Msg} service and executes only as the payload of a
threshold-approved coordinator action
(\texttt{vote-sdk/x/vote/keeper/msg\_server\_coordinator\_actions.go:294-298}).
The handler derives the round identifier on-chain as one Poseidon hash
over the creation height, the snapshot block hash and the proposals hash
(each split into two field elements), the end time, and both anchors, so
\texttt{vote\_round\_id} commits to the snapshot
(\texttt{tx.proto:32-35}; \texttt{msg\_server.go:30-48}).

\paragraph{Message 2: \texttt{MsgDelegateVote} (ZKP\,1).} The delegation
message carries \texttt{rk} (the randomized spend-validating key),
\texttt{spend\_auth\_sig} (RedPallas), \texttt{signed\_note\_nullifier} (the
nullifier of the synthetic keystone note), \texttt{cmx\_new} (the output
note commitment), \texttt{van\_cmx} (the Vote Authority Note commitment),
\texttt{gov\_nullifiers} (up to five governance nullifiers),
\texttt{proof}, \texttt{vote\_round\_id}, and \texttt{sighash} --- a
\emph{client-provided} 32-byte value over which the chain verifies the
RedPallas signature without recomputing it (\texttt{tx.proto:55-65}). This
section proves the eligibility half of the message; the Vote Authority Note
is constructed in \S\ref{sec:vote-van}, and the synthetic keystone note
with its signature binding --- the $\rho$-chain through which the RedPallas
signature commits to the whole delegation --- in
\S\ref{sec:vote-keystone}.

\paragraph{Message 3: \texttt{MsgCastVote} (ZKP\,2).} A cast vote carries
\texttt{van\_nullifier},
\texttt{vote\_\allowbreak authority\_\allowbreak note\_\allowbreak new}
(the successor
Vote Authority Note), \texttt{vote\_commitment}, \texttt{proposal\_id},
\texttt{proof}, \texttt{vote\_round\_id}, the anchor height
\texttt{vote\_\allowbreak comm\_\allowbreak tree\_\allowbreak
anchor\_\allowbreak height}, \texttt{vote\_auth\_sig}, and
\texttt{r\_vpk}, a compressed Pallas point that the chain decompresses for
the verifier; unlike message~2, the sighash here is recomputed on-chain
(\texttt{tx.proto:72-84}, including the removed-field comments).

\paragraph{Message 4: \texttt{MsgRevealShare} (ZKP\,3).} A share reveal
carries \texttt{share\_nullifier}, \texttt{enc\_share} --- 64 bytes, one
lifted-ElGamal ciphertext $C_1 \,\|\, C_2$ as compressed Pallas points ---
\texttt{proposal\_id}, \texttt{vote\_decision}, \texttt{proof},
\texttt{vote\_round\_id}, and \texttt{vote\_comm\_tree\_anchor\_height}
(\texttt{tx.proto:89-97}). Helper servers, not wallets, submit it: the
share nullifier is constructed in \S\ref{sec:vote-layers}, and the helper
protocol and the trust it grants are treated in
\S\ref{sec:vote-trust-map}.

\paragraph{Message 5: \texttt{MsgSubmitTally}.} The closing message carries
\texttt{vote\_round\_id}, \texttt{creator} --- which the chain checks
against the \emph{block proposer}: the proposer's node injects the message
during tallying, and the proto comment naming the session creator is
unenforced (\texttt{x/vote/keeper/msg\_server\_tally\_decrypt.go:77-79};
\texttt{app/prepare\_proposal.go:194-212}) --- and one \texttt{TallyEntry}
$\{$\texttt{proposal\_id}, \texttt{vote\_decision}, \texttt{total\_value},
\texttt{decryption\_proof}$\}$ per (proposal, decision) pair
(\texttt{tx.proto:104-116}). The \texttt{total\_value} field is denominated
in ballots despite its proto comment saying zatoshi; we return to this in
\S\ref{sec:vote-decrypt}.

\paragraph{Classification.} Designed but unspecified: the message set,
field semantics, and transport exist only in
\texttt{vote-sdk/proto/svote/v1/tx.proto} at commit \texttt{cb915f51}
(Table~\ref{tab:vote-sources}). A specification would need to pin the
byte-level encoding and length of every field, the validation order, the
envelope-bypass transport, and the round-identifier derivation --- the last
currently fixed only by a proto comment and the FFI implementation it
describes.

\subsection{Eligibility: the note-bundle relation (ZKP\,1)}
\label{sec:vote-zkp1}

To delegate weight, a holder must prove three things about a set of Orchard
notes without revealing any of them: the notes are theirs; each existed at
the snapshot; each was unspent at the snapshot. The delegation circuit
proves the conjunction --- one Halo\,2/IPA circuit over Pallas at $K=14$
($16{,}384$ rows) with $14$ public inputs
(\texttt{voting-circuits/src/delegation/circuit.rs:111} and the module
README, both at commit \texttt{4c39abd5}).

\paragraph{Bundle shape.} The circuit carries exactly five note slots
(\texttt{MAX\_REAL\_NOTES}); a wallet with fewer notes pads the remaining
slots with zero-value notes, and one with more produces multiple delegation
proofs. The fixed width is deliberate: every delegation publishes the same
shape --- five governance nullifiers at public-input offsets $8$--$12$ ---
so the published proof does not reveal how many real notes it bundles
(\texttt{src/delegation/README.md}, Inputs). The two ledger anchors enter as
public inputs \texttt{nc\_root} (offset $6$) and \texttt{nf\_imt\_root}
(offset $7$); the nullifier domain $\mathsf{dom}$ (offset $13$) is exposed
for API compatibility but constrained in-circuit to
$\mathsf{Poseidon}(\mathsf{tag}, \texttt{vote\_round\_id})$, where
$\mathsf{tag}$ is the ASCII string \texttt{governance authorization}
zero-padded to a Pallas base-field element
(\texttt{src/domain\_tags.rs}; \texttt{src/delegation/imt.rs:26-35}).
Poseidon throughout the protocol is the P128Pow5T3 instance, width $3$,
rate $2$ (\texttt{src/protocol\_hash.rs}).

\paragraph{Per-slot conditions.} Each slot proves six conditions, numbered
$9$--$14$ in the source\footnote{The headline counts disagree ---
\texttt{circuit.rs} line~3 says ``all 14 conditions'' but line~108 ``all 15
conditions''; \texttt{README.md} lines~3 and~52 say ``15'' and ``conditions
9--15'' --- yet both files enumerate exactly conditions 1--14 and neither
defines a condition~15: the README's numbered sections run 1 through 14 and
the per-slot list is 9--14 throughout. We state the total as fourteen ---
eight keystone-global conditions plus six per-slot conditions instantiated
five times --- and record the ``15'' as a stale count in the sources.}
(\texttt{src/delegation/circuit.rs:1497-1805}):
\begin{itemize}
\item \emph{Note commitment integrity} (condition~9). The prover witnesses
  the note plaintext $(\gd, \pkd, v, \rho, \psi)$, the trapdoor $\rcm$, and
  the claimed commitment $\cmsym$; the circuit recomputes
  $\NoteCommit_{\rcm}$ over the witnessed fields, constrains it equal to
  $\cmsym$, and extracts $\cmx = \Extract(\cmsym)$ as an internal wire. The
  commitment scheme is the deployed Orchard \NoteCommit{} --- specified
  bin, reused unchanged through the \texttt{unstable-voting-circuits}
  feature (\S\ref{sec:vote-source-situation}).
\item \emph{Address ownership} (condition~11). The circuit checks
  $\pkd = [\ivk^{*}]\,\gd$, where $\ivk^{*}$ is selected by a witnessed
  boolean $\mathsf{is\_internal}$ through the custom gate
  \texttt{q\_scope\_select},
  $\ivk^{*} = \ivk + \mathsf{is\_internal}\cdot(\ivk_{\mathrm{int}} - \ivk)$;
  both candidates are derived in-circuit by \CommitIvk{} from the witnessed
  $\rivk$ and its internal-scope counterpart (condition~5, global). Change
  notes therefore carry weight on equal terms with externally received
  ones. A wrong scope flag selects the wrong key and the ownership check
  rejects.
\item \emph{Snapshot membership} (condition~10). The circuit computes the
  \Sinsemilla{} Merkle root from the leaf $\cmx$ over the witnessed
  $32$-level authentication path, and the per-note gate
  \texttt{q\_per\_note} enforces
  \[
    v \cdot (\mathsf{root} - \mathsf{anchor}) = 0,
  \]
  with $\mathsf{anchor}$ the public \texttt{nc\_root}
  (\texttt{circuit.rs:1777-1798}). A real note ($v \neq 0$) must resolve to
  the anchor; a zero-value padding note --- which has no leaf in the Zcash
  tree --- satisfies the gate vacuously.
\item \emph{Real nullifier} (condition~12). The circuit derives the note's
  true Orchard nullifier
  $\nfsym = \mathsf{DeriveNullifier}_{\nksym}(\rho, \psi, \cmsym)$ --- the
  deployed Orchard rule, constructed in the \emph{Ironwood Guide}
  (\S``Prevention of double-spending: nullifiers'') --- in-circuit and
  \emph{never publishes it}: it exists only as a private wire feeding
  conditions~13 and~14 (\texttt{circuit.rs:1682-1701}).
\item \emph{Snapshot non-spentness} (condition~13). The circuit proves
  non-membership of $\nfsym$ in the snapshot's nullifier set under the
  public \texttt{nf\_imt\_root} --- the construction of
  \S\ref{sec:vote-imt}. The root equality is unconditional: padding notes,
  too, must carry a valid non-membership proof
  (\texttt{README.md} \S13).
\item \emph{Governance nullifier} (condition~14). The circuit computes and
  publishes
  $\nfsym^{\mathsf{gov}} = \mathsf{Poseidon}(\nksym, \mathsf{dom}, \nfsym)$
  at the slot's public-input offset. Keyed by the private $\nksym$, the
  published value stays unlinkable to $\nfsym$ even after the note is later
  spent on the main chain (\texttt{circuit.rs:1703-1729}). Its role in
  double-vote prevention is treated in \S\ref{sec:vote-layers}. Chain-side,
  the delegation handler applies the same atomic check-and-set to all five
  published slots, padding included
  (\texttt{vote-sdk/x/vote/keeper/msg\_server.go:148-154}); a padding
  slot's nullifier enters the same round-scoped set, unique because the
  wallet samples fresh $\rho$ and \texttt{rseed} per synthetic padding
  note (\texttt{src/delegation/builder.rs}, \texttt{build\_padding\_slot}),
  and a re-used padding witness is rejected as a duplicate delegation ---
  a completeness cost, never a soundness one.
\end{itemize}

The conjunction across a slot reads: the prover's key material owns a note
whose commitment lies in the Zcash note-commitment tree at the snapshot
(membership under \texttt{nc\_root}) and whose nullifier had not been
revealed by the snapshot (non-membership under \texttt{nf\_imt\_root}) ---
existed, and unspent, at the pinned height.

\paragraph{Classification.} The relation as a whole is designed but
unspecified: it exists only as the circuit at
\texttt{voting-circuits}~@~\texttt{4c39abd5}, which is simultaneously the
prover and the verifier (\S\ref{sec:vote-source-situation}). Its
ingredients \NoteCommit{}, nullifier derivation, \CommitIvk{}, and the
\Sinsemilla{} Merkle instance are the specified bin, inherited through
zcash/orchard~@~\texttt{c0b6b917}. A specification would need to pin the
public-input layout, the custom-gate semantics, the fixed five-slot shape,
and the padding-note rules.

\subsection{Snapshot non-spentness: the Poseidon indexed Merkle tree}
\label{sec:vote-imt}

Membership alone would let a spent note vote: the note-commitment tree is
append-only and a spent note's commitment remains in it forever. Snapshot
eligibility therefore needs a second, negative statement --- the note's
nullifier is \emph{not} in the set of nullifiers revealed up to the snapshot
height. Proving non-membership in a Merkle set requires structure beyond a
plain accumulator; the protocol uses an indexed Merkle tree over the sorted
nullifier set.

\paragraph{Construction.} The tree is Poseidon-based with depth $29$
(\texttt{IMT\_DEPTH},
\texttt{voting-circuits/\allowbreak src/\allowbreak
delegation/\allowbreak imt.rs:19}).
Each leaf is a \emph{punctured range} --- a triple of sorted boundaries,
hashed and covering the punctured open interval
\[
\mathsf{leaf} =
\mathsf{Poseidon}(\nfsym_{\mathrm{lo}}, \nfsym_{\mathrm{mid}},
\nfsym_{\mathrm{hi}}),
\qquad
(\nfsym_{\mathrm{lo}}, \nfsym_{\mathrm{hi}}) \setminus
\{\nfsym_{\mathrm{mid}}\},
\]
so the set members appear only as interval boundaries or interior
punctures, and everything strictly between them is certified absent.
In-circuit, non-membership of the private $\nfsym$ costs
$31$ Poseidon calls ($2$ for the leaf, $29$ for the path) and proves
(\texttt{imt.rs:51-70}; \texttt{src/delegation/README.md} \S13):
\begin{enumerate}
\item the leaf authenticates: the $29$-level path from the leaf hash
  resolves to the public \texttt{nf\_imt\_root};
\item the nullifier lies strictly inside the interval: the offsets
  $x_{\mathrm{lo}} = \nfsym - \nfsym_{\mathrm{lo}} - 1$ and
  $x_{\mathrm{hi}} = \nfsym_{\mathrm{hi}} - \nfsym - 1$ are each
  range-checked to $[0, 2^{250})$; and
\item the nullifier misses the puncture:
  $(\nfsym - \nfsym_{\mathrm{mid}}) \cdot \delta = 1$ for a witnessed
  inverse $\delta$.
\end{enumerate}

\paragraph{What it excludes --- and what it does not.} The statement
excludes exactly the nullifiers revealed on the Zcash chain up to the
snapshot height: the note was unspent when the snapshot was taken. It
excludes nothing after that height --- a note spent the block after the
snapshot still delegates its weight, because snapshot voting fixes the
electorate at a single height by design. Two soundness assumptions live
\emph{outside} the circuit (\texttt{imt.rs:56-70};
\texttt{README.md} \S13). First, the tree contract: the circuit checks
only the one authenticated leaf, so the committed leaves must form the
canonical punctured-range tree for the set --- sorted, deduplicated,
non-overlapping except for shared boundaries. A malformed tree in which
some set member lies strictly inside another leaf's interval would certify
a spent note as unspent, and no per-leaf gate can detect the global
overlap. Second, the span bound: the $250$-bit offset checks require
$\nfsym_{\mathrm{hi}} - \nfsym_{\mathrm{lo}} \le 2^{250}$ in canonical
field order, which the builder guarantees by pre-populating $33$ sentinel
nullifiers spaced at most $2^{249}$ apart, plus $p-1$ to close the tail
(the \texttt{SpacedLeafImtProvider} strategy). The tree is built by the
PIR operator and its root posted by the coordinator --- which is where the
next subsection picks up.

\paragraph{Classification.} Designed but unspecified
(\texttt{voting-circuits}~@~\texttt{4c39abd5}). A specification would need
to pin the leaf model, the sentinel schedule, and above all the
canonical-tree contract --- including which party is accountable for the
tree's well-formedness, since the circuit cannot check it.

\subsection{The snapshot roots: posted, never verified}
\label{sec:vote-snapshot}

Both anchors originate off-chain, in coordinator tooling. The
note-commitment root \texttt{nc\_root} is recomputed by Rust FFI as the
\Sinsemilla{} hash of an Orchard frontier fetched from a public
lightwalletd server; the nullifier root \texttt{nullifier\_imt\_root} is
fetched from the running PIR service
(\texttt{vote-sdk/api/snapshot.go:44-66}). The coordinator posts both in
message~1; the handler stores them verbatim in the round record and folds
them into \texttt{vote\_round\_id}
(\texttt{x/vote/keeper/msg\_server.go:36-138}); and at validation the ante
handler looks the stored roots up from the round and passes them to the
FFI verifier as the public inputs of ZKP\,1
(\texttt{x/vote/ante/validate.go:147-177}).

Nothing on the vote chain verifies that the posted roots match real Zcash
state: \texttt{svoted} runs no Zcash light client, and the circuit
documentation is explicit that it ``proves statements relative to the
public inputs supplied by the verifier; it does not authenticate where
those inputs came from''
(\texttt{voting-circuits/src/delegation/README.md}, Public Input
Provenance). The trust structure that results is asymmetric. Honest
wallets are safe from a \emph{wrong} snapshot in the sense that their
proofs simply fail: their witnesses --- Merkle paths and IMT leaves ---
come from the real chain and cannot resolve to fabricated roots. And the
roots are public, so anyone with a Zcash node can audit them after the
fact. But nothing prevents a coordinator quorum from pinning fabricated
roots for which it alone holds matching witnesses, minting arbitrary
voting weight that the chain would accept; detection is post hoc and
protocol-external. The trust root for eligibility is therefore the
coordinator threshold together with the PIR operator that builds the
nullifier tree of \S\ref{sec:vote-imt}.

\paragraph{Classification.} Designed but unspecified, with an explicit
trust assumption that must be flagged wherever eligibility is discussed. A
specification would need either chain-side verification (a light client or
a succinct proof that the roots match a Zcash block) or, failing that, a
normative audit procedure with a defined challenge window.

\subsection{Weight: ballot quantization}\label{sec:vote-weight}

Delegated weight is denominated in ballots. The circuit sums the five slot
values into $v_{\mathrm{total}} = \sum_{i=1}^{5} v_i$ (internal wires
produced by condition~9, consumed by conditions~7 and~8;
\texttt{src/delegation/README.md}, Inputs), and one ballot represents
$D = 12{,}500{,}000$ zatoshi ($0.125$~ZEC): the constant
\texttt{BALLOT\_DIVISOR} (\texttt{voting-circuits/src/params.rs:14}). The
honest ballot count is $q = \lfloor v_{\mathrm{total}} / D \rfloor$.

Condition~8 (\texttt{src/delegation/circuit.rs:1290-1427}) witnesses the
count $q$ (\texttt{num\_ballots}) and a remainder $r$ as \emph{free
advice}, and constrains
\begin{equation}\label{eq:vote-quant}
q \cdot D + r = v_{\mathrm{total}}, \qquad
0 \le r < 2^{24}, \qquad
1 \le q \le 2^{30}.
\end{equation}
The range checks run through a $10$-bit lookup table: the circuit
multiplies $r$ by $2^{6}$ and checks the product to $30$ bits (so
$r < 2^{24}$), and checks $q - 1$ to $30$ bits directly --- which enforces
both bounds at once, since $q = 0$ wraps $q-1$ to $p - 1 \approx 2^{254}$
and fails. The lower bound $q \ge 1$ doubles as a minimum stake: a bundle
worth less than $0.125$~ZEC admits no valid witness. The upper bound is
safe slack: $2^{30}$ ballots is about $134$ million ZEC, far above the
$21$ million supply. The remainder bound $2^{24}$ is the tightest
power-of-two-aligned bound that keeps completeness --- every honest
$r < D$ fits --- but it is \emph{looser} than $D$, and that slack admits a
second witness.

\begin{proposition}[One-ballot under-claim window]\label{thm:vote-underclaim}
Fix $v_{\mathrm{total}}$ and let
$\hat{q} = \lfloor v_{\mathrm{total}}/D \rfloor$ and
$\hat{r} = v_{\mathrm{total}} \bmod D$ form the honest witness, assumed
valid ($\hat{q} \ge 1$). The pairs satisfying the constraints of
Equation~\eqref{eq:vote-quant} are exactly $(\hat{q}, \hat{r})$ together
with $(\hat{q} - 1,\; \hat{r} + D)$, the latter admissible if and only if
\[
\hat{r} < 2^{24} - D = 4{,}277{,}216
\qquad\text{and}\qquad
\hat{q} \ge 2 .
\]
In particular, a prover can claim exactly one ballot fewer than floor
division yields --- for approximately $34.2\%$ of totals under a uniform
residue model, since $(2^{24} - D)/D \approx 0.3422$ --- and can never
claim more.
\end{proposition}

\begin{proof}
Suppose $(q_1, r_1)$ and $(q_2, r_2)$ both satisfy
Equation~\eqref{eq:vote-quant} for the same $v_{\mathrm{total}}$.
Subtracting the two instances of the linear constraint gives
$(q_1 - q_2)\,D \equiv r_2 - r_1 \pmod{p}$. The range checks bound the
canonical representatives: $|q_1 - q_2| \cdot D \le (2^{30} - 1)D <
2^{54}$ and $|r_2 - r_1| < 2^{24}$, both far below $p > 2^{254}$, so the
congruence holds over the integers. Then
$|q_1 - q_2| = |r_2 - r_1| / D < 2^{24}/D < 2$, forcing
$|q_1 - q_2| \in \{0, 1\}$. The case $0$ gives $r_1 = r_2$. The case
$q_1 = q_2 + 1$ gives $r_2 = r_1 + D$, which is admissible precisely when
$r_2 < 2^{24}$ (the remainder range check) and $q_2 \ge 1$ (the count
range check) --- for the honest pair, the stated window. The over-claim
direction from the honest witness, $(\hat{q} + 1, \hat{r} - D)$, requires
$\hat{r} - D < 0$; its canonical field representative is at least
$p - D$, which fails the shifted $30$-bit range check. No further pairs
exist.
\end{proof}

\paragraph{Impact and status.} The deviation is self-harming: the only
dishonest witness \emph{under}-claims, discarding one ballot of the
prover's own weight, and aggregate totals consequently err low, never
high. No downstream check repairs or re-derives the count: the Vote
Authority Note commitment (condition~7) hashes whatever $q$ the prover
chose, and the vote-proof circuit later opens the same value. The gap is
documented by the implementers themselves ---
\texttt{src/delegation/README.md} \S8, ``Soundness scope'', carries the
analysis above together with two single-constraint tightenings (a
$24$-bit check on $(D - 1) - r$, or on $r + (2^{24} - D)$, either of
which pins $r \in [0, D)$ at a cost of roughly $3$--$4$ rows) --- and it
is unfixed at commit \texttt{4c39abd5}. The public documentation states
exact floor division (gitbook ZKP\,1 page, fetched 2026-07-08), which the
circuit does not enforce: this is the single spec-versus-code
disagreement recorded in this volume
(\S\ref{sec:vote-source-situation}), conservative in direction --- the
code is the protocol, and the code is weaker.

\paragraph{Classification.} The quantization relation as implemented ---
Equation~\eqref{eq:vote-quant} --- is designed but unspecified. The gap
between it and the stated floor-division rule is this volume's one
\emph{open problem}: the design intends exact division, the circuit
proves less, and the tightening exists on paper only. A specification
must either adopt the weaker relation and its under-claim window
explicitly, or mandate one of the documented tightening constraints.

\section{Vote encryption, nullifiers, and tally}\label{sec:vote-tally}

Delegation leaves the governance hotkey holding a Vote Authority Note (VAN): a
blinded commitment to the hotkey address, a ballot count, and a per-proposal
authority mask. Three problems remain. The hotkey must cast a vote whose
weight never appears in public; the protocol must stop every form of
double-voting --- the same coins delegated twice, the same VAN spent twice,
the same proposal voted twice along a chain of successor VANs, the same
encrypted share counted twice; and the aggregate totals must become public in
a form every chain node verifies, without any single party ever holding the
decryption key. This section develops the answers in order: the VAN itself,
the synthetic keystone note whose signature authorises its minting, a
four-layer double-vote discipline, a sixteen-share lifted-ElGamal encoding, a
per-round threshold Election Authority, and a chain-verified tally.

\paragraph{Classification.} Every governance-specific object in this section
is \emph{designed but unspecified}: the normative artifact is the
implementation, cited commit-pinned ---
\texttt{valargroup/voting-circuits} at \texttt{4c39abd5} (2026-06-03) and
\texttt{valargroup/vote-sdk} at \texttt{cb915f51} (2026-06-05). File paths
beginning \texttt{src/} refer to the former; paths beginning
\texttt{x/vote/}, \texttt{crypto/}, or \texttt{proto/} to the latter. The
would-be specification (zcash/zips PR~\#1200, opened 2026-03-05) remains
unmerged as of 2026-07-08, and the single \texttt{voting-circuits} crate both
proves and verifies, so no independent implementation cross-checks these
formulas. The \emph{specified} bin holds only the inherited Orchard
primitives --- \NoteCommit, nullifier derivation, \CommitIvk, RedPallas, the
Sinsemilla Merkle instance --- reused through the \texttt{orchard} crate's
\texttt{unstable-voting-circuits} feature (zcash/orchard commit
\texttt{c0b6b917}). One \emph{open problem} reaches into this section from
the delegation circuit; Remark~\ref{rem:vote-quant} flags it where it lands.

Throughout this section, $\mathsf{Poseidon}$ denotes the P128Pow5T3 instance
over $\F_{p_{\Pallas}}$ (width $3$, rate $2$), applied with the fixed-length
padding of the stated arity. Domain separation uses two encodings, both
pinned by unit tests (\texttt{src/domain\_tags.rs}): small numeric tags ($0$
for VAN commitments, $1$ for vote commitments), and short ASCII tags ---
\texttt{"governance authorization"}, \texttt{"vote authority spend"},
\texttt{"share spend"} --- zero-padded to $32$ bytes and read as little-endian
integers, which are canonical in $\F_{p_{\Pallas}}$ because the top byte
stays zero. We abbreviate the ASCII tags $T_{\mathsf{gov}}$,
$T_{\mathsf{spend}}$, $T_{\mathsf{share}}$, and write $\mathsf{rid}$ for the
voting round identifier and $\mathsf{pid}$ for the proposal identifier.

\subsection{The Vote Authority Note}\label{sec:vote-van}

The VAN is the vote chain's carrier of authority: one field element that
commits to who may vote (the hotkey address), with how much weight (the
ballot count), in which round, and on which proposals (the authority mask).

\begin{definition}[VAN commitment]\label{def:vote-van}
Let $(\gd, \pkd)$ be the hotkey's diversified address in the voting key
hierarchy, $n_b$ the ballot count, $\mathsf{auth} \in [0, 2^{16})$ the
proposal-authority mask, and $\mathsf{r}_{\mathsf{van}}$ a uniformly random
blind. The VAN commitment is the two-layer hash
\[
\mathsf{van} := \mathsf{Poseidon}\bigl(\mathsf{core},\
\mathsf{r}_{\mathsf{van}}\bigr), \qquad
\mathsf{core} := \mathsf{Poseidon}\bigl(0,\ x(\gd),\ x(\pkd),\ n_b,\
\mathsf{rid},\ \mathsf{auth}\bigr),
\]
with $x(\cdot)$ the $x$-coordinate
(\texttt{src/gadgets/van\_integrity.rs}, lines 61--78; one gadget is shared
by the delegation and vote-proof circuits).
\end{definition}

The two layers divide labour. The inner layer is structural, and its preimage
is low-entropy: the round id is public, the mask at mint is a fixed constant,
and ballot counts are small integers --- a bare inner hash would admit
dictionary attacks on the weight and the address. The outer layer closes them
with the uniform blind (module documentation,
\texttt{src/gadgets/van\_integrity.rs}, lines 23--27).

The address enters by $x$-coordinates alone, which identifies
$(\gd, \pkd)$ only up to coordinated negation. The construction is safe
because delegation additionally binds the full points through \NoteCommit{}
into the public $\cmx_{\mathrm{new}}$; the module documentation warns that
any future caller dropping that side constraint must widen the preimage to
full points (\texttt{src/gadgets/van\_integrity.rs}, lines 14--22).

The mask $\mathsf{auth}$ is a $16$-bit bitmask minted at
$2^{16} - 1 = 65535$ --- all bits set, a constant baked into the verification
key (\texttt{src/delegation/circuit.rs}, lines 162 and 1464--1467). Bit $0$
is a
permanent sentinel that no vote may consume, so usable bits are $1$ through
$15$ and a round supports at most $15$ proposals
(\texttt{src/vote\_proof/circuit.rs}, lines 123--143).

The vote chain keeps one per-round Poseidon Merkle tree of depth $24$
(\texttt{src/params.rs}, line 25; per-round state in
\texttt{x/vote/keeper/keeper\_voting.go}, lines 72--80): delegation appends
the freshly minted VAN commitment; each cast appends the successor VAN and
then the vote commitment. The reduced depth is a sizing choice ---
governance produces one leaf per delegation plus two per vote, well within
$2^{24} \approx 1.7 \times 10^{7}$ leaves (\texttt{src/params.rs}, lines
16--25).

\subsection{Keystone authorisation: the synthetic signed note}
\label{sec:vote-keystone}

The delegation message is authorised by a RedPallas spend-auth signature
under the randomized key $\rk$ (\S\ref{sec:vote-messages}, message~2), and
the intended signer is a Keystone-class hardware wallet, which signs only
Orchard Actions. The voting client therefore wraps the delegation in an
Orchard-Action shape around a \emph{synthetic} keystone note: a
$1$-zatoshi note that exists on no chain --- the circuit witnesses no
Merkle path for it and checks no anchor; the value is $1$ rather than $0$
only because the wallet UI does not render zero-value spends
(\texttt{src/delegation/README.md}, ``Integration: the Keystone (signed)
note is synthetic'').

Conditions~1--6 of the delegation circuit govern this branch
(\texttt{src/delegation/circuit.rs}, lines 16--29); the load-bearing pair
is conditions~3 and~2. Condition~3 forces the keystone note's $\rho$ to a
Poseidon hash of everything the delegation publishes,
\[
\rho_{\mathrm{signed}} = \mathsf{Poseidon}\bigl(\cmx_1, \ldots, \cmx_5,\
\mathsf{van},\ \mathsf{rid}\bigr),
\]
the arity-$7$ fixed-length instance (\texttt{rho\_binding\_hash},
\texttt{circuit.rs:176-190}). Condition~2 then derives the keystone
nullifier by the deployed Orchard rule --- constructed in the
\emph{Ironwood Guide} (\S``Prevention of double-spending: nullifiers'') ---
\[
\nfsym_{\mathrm{signed}} =
\mathsf{DeriveNullifier}_{\nksym}\bigl(\rho_{\mathrm{signed}},\
\psi_{\mathrm{signed}},\ \cmsym_{\mathrm{signed}}\bigr),
\]
and exposes it as a public input. The chain verifies the RedPallas
signature under $\rk$ over the \emph{client-provided} sighash and checks
that the proof's $\nfsym_{\mathrm{signed}}$ equals the wrapping Action's
nullifier (\texttt{x/vote/ante/validate.go}), so the hardware wallet's
signature commits --- through the chain
$\rk \rightarrow \nfsym_{\mathrm{signed}} \rightarrow
\rho_{\mathrm{signed}} \rightarrow \mathsf{van}$ --- to the five bundle
commitments, the VAN, and the round identifier. Removing any link breaks
the binding; this chain is why the chain-side acceptance of a
client-provided sighash is safe. The remaining conditions are supporting:
condition~4 proves $\rk = [\alpha]\,\mathsf{SpendAuthG} + \aksym$ for a
witnessed randomizer $\alpha$ (the \emph{Ironwood Guide},
\S``Authorisation of the spend: RedPallas''); condition~5 ties $\aksym$,
$\nksym$, and $\rivk$ to the same \CommitIvk{} chain the real-note slots
use (\S\ref{sec:vote-zkp1}, condition~11); conditions~1 and~6 recompute
the keystone and output note commitments from witnessed openings ---
knowledge checks, not membership checks
(\texttt{src/delegation/README.md}, the three-bucket audit note).

\subsection{Double-vote prevention: four layers}\label{sec:vote-layers}

A ballot can be replayed at four distinct joints: the Zcash note (delegate it
twice), the VAN (spend it twice), the proposal bit (vote the same proposal
again through a successor VAN), and the encrypted share (feed it to the
accumulator twice). One defence guards each joint. The first, second, and
fourth are nullifiers; chain-side, each kind lands in its own type- and
round-scoped set under an atomic check-and-set
(\texttt{x/vote/keeper/keeper\_voting.go}, lines 19--48; type constants
\texttt{x/vote/types/keys.go}, lines 96--101). The third is arithmetic,
enforced in-circuit.

\paragraph{Layer 1: the governance nullifier (delegation).} Each delegated
note exposes
\[
\nfsym^{\mathsf{gov}} := \mathsf{Poseidon}\bigl(\nksym,\ \mathsf{dom},\
\nfsym\bigr), \qquad
\mathsf{dom} := \mathsf{Poseidon}\bigl(T_{\mathsf{gov}},\ \mathsf{rid}\bigr),
\]
where $\nksym$ and $\nfsym$ are the \emph{Zcash} note's nullifier key and
real Orchard nullifier, both computed in-circuit --- the real nullifier is
never published (\texttt{src/delegation/circuit.rs}, lines 1703--1729). The
domain $\mathsf{dom}$ is a public input, but the circuit re-derives it from
the public round id and constrains equality
(\texttt{src/delegation/circuit.rs}, lines 1041--1055), so a prover cannot
substitute a foreign domain to escape the round scoping. The chain rejects
duplicates in the governance-typed, round-scoped set. The attack closed:
delegating the same unspent-at-snapshot note into two bundles in one round,
which would double its weight. Keying by the secret $\nksym$ serves privacy:
when the note is later spent on mainnet and its real nullifier becomes
public, no observer can link the pair --- the governance trail does not
deanonymise the payment trail.

\paragraph{Layer 2: the VAN nullifier (casting).} Spending a VAN publishes
\[
\nfsym^{\mathsf{van}} := \mathsf{Poseidon}\bigl(\nksym_{v},\
T_{\mathsf{spend}},\ \mathsf{rid},\ \mathsf{van}_{\mathrm{old}}\bigr),
\]
where $\nksym_{v}$ is the voting hierarchy's nullifier-deriving key, proved
to belong to the witnessed address through the \CommitIvk{} chain
(\texttt{src/vote\_proof/circuit.rs}, lines 203--222). A cast proves
membership of $\mathsf{van}_{\mathrm{old}}$ in the round tree without
identifying the leaf, and this nullifier is the only spend mark. The attack
closed: re-spending the old VAN, whose mask still carries the bit the vote
just consumed --- without the nullifier, one delegation would vote the same
proposal arbitrarily often from the same leaf.

\paragraph{Layer 3: the authority-mask decrement (casting).} Layer~2 kills
the parent, not the lineage: every cast mints a successor VAN, and a
successor still carrying the consumed bit would re-enable the same vote.
Condition~6 of the vote-proof circuit therefore forces
\[
\mathsf{auth}_{\mathrm{new}} = \mathsf{auth}_{\mathrm{old}} -
2^{\mathsf{pid}}.
\]
Field arithmetic cannot express a variable exponent as a polynomial gate, so
the prover witnesses $2^{\mathsf{pid}}$ and a lookup over the table
$\{(j, 2^{j})\}_{j=0}^{15}$ proves it; a field-inverse gate excludes
$\mathsf{pid} = 0$ (the sentinel); a $16$-bit decomposition of
$\mathsf{auth}_{\mathrm{old}}$ with a one-hot selector anchored at
$\mathsf{pid}$ proves the consumed bit was actually set ($\sum_j
\mathsf{sel}_j b_j = 1$), and the recomposition returns
$\mathsf{auth}_{\mathrm{new}}$ with exactly that bit cleared, implicitly
range-checked to $16$ bits
(\texttt{src/vote\_proof/gadgets/authority\_decrement.rs}, lines 34--93).
The successor VAN commits $\mathsf{auth}_{\mathrm{new}}$ and re-enters the
tree. The invariant bought by layers 2 and 3 together: one delegation yields
at most one vote per proposal, and at most $15$ proposals per round.

\paragraph{Layer 4: the share nullifier (reveal).} Revealing the $i$-th
encrypted share of a vote publishes
\[
\nfsym^{\mathsf{share}}_{i} := \mathsf{Poseidon}\bigl(T_{\mathsf{share}},\
\mathsf{vc},\ i,\ \mathsf{blind}_i\bigr),
\]
\begin{sloppypar}
where $\mathsf{vc}$ is the vote commitment (\S\ref{sec:vote-elgamal}) and
$\mathsf{blind}_i$ the share-commitment blind
(\texttt{src/share\_reveal/circuit.rs}, lines 161--187). The attack closed:
the tally accumulator is a raw homomorphic sum
(\S\ref{sec:vote-decrypt}), so re-submitting one accepted ciphertext would
double-count its weight; the share-typed set admits each
$(\mathsf{vc}, i)$ once. The blind is load-bearing for privacy rather than
soundness: the vote commitment is public in the cast message and $i$ ranges
over sixteen values, so an unblinded nullifier could be evaluated by any
observer against every commitment in the tree, linking each reveal to its
vote by dictionary test --- collapsing exactly the anonymity that the reveal
proof's non-identifying membership provides. Blinds never appear on chain.
\end{sloppypar}

\subsection{Vote encoding: sixteen lifted-ElGamal shares}
\label{sec:vote-elgamal}

Casting must convey weight to the tally without publishing it per vote. The
instrument is additively homomorphic encryption to a key that no single
party holds (\S\ref{sec:vote-dkg}).

\begin{definition}[Vote share encryption]\label{def:vote-elgamal}
Fix the generator
\[
G := \GroupHash^{\Pallas}(\mathtt{z.cash:Orchard}, \mathtt{G}),
\]
Orchard's spend-authorisation base. A vote on proposal
$\mathsf{pid}$ splits its ballot count $n_b$ into sixteen prover-chosen
shares $v_0, \ldots, v_{15}$ subject to
\[
\sum_{i=0}^{15} v_i = n_b, \qquad v_i \in [0, 2^{30}),
\]
and encrypts each under the round's Election Authority key
$\mathsf{ea_{pk}} \in \Ep(\F_{p_{\Pallas}})$:
\[
C_{1,i} := [r_i]\,G, \qquad
C_{2,i} := [v_i]\,G + [r_i]\,\mathsf{ea_{pk}},
\]
with fresh nonzero randomness $r_i$
(\texttt{src/gadgets/elgamal.rs}; sum and range are conditions 8 and 9,
\texttt{src/vote\_proof/circuit.rs}, lines 416--457 and 1026--1036).
\end{definition}

The scheme is \emph{lifted}: the message enters the exponent, so
componentwise addition of ciphertexts encrypts the sum of plaintexts and the
accumulator needs no key. The price is paid at decryption, which yields
$[v]\,G$ and recovers $v$ only as a bounded discrete logarithm --- the
$2^{28}$ bound of \S\ref{sec:vote-decrypt}. Reusing $\mathsf{SpendAuthG}$
avoids introducing a second fixed base; the in-circuit $22$-window
short-scalar tables share the same generator
(\texttt{src/gadgets/elgamal.rs}, lines 57--66). The range check closes
wraparound: shares are field elements, and without it a decomposition could
encode negative shares as large residues while still satisfying the sum
constraint. Sixteen shares serve weight privacy under staggered reveal: the
prover-chosen random decomposition leaves no weight fingerprint on any
single revealed ciphertext (\texttt{src/vote\_proof/circuit.rs}, lines
418--423). Client-side, the randomness $r_i$ and the blinds derive from a
BLAKE2b-512 PRF with personalisation \texttt{ZcashVote\_Expand} under
distinct domain bytes (\texttt{src/domain\_tags.rs}, lines 28--41).

The cast binds everything into one public field element:
\[
\mathsf{vc} := \mathsf{Poseidon}\bigl(1,\ \mathsf{rid},\
h_{\mathsf{sh}},\ \mathsf{pid},\ d\bigr), \qquad
h_{\mathsf{sh}} := \mathsf{Poseidon}\bigl(\mathsf{sc}_0, \ldots,
\mathsf{sc}_{15}\bigr),
\]
\[
\mathsf{sc}_i := \mathsf{Poseidon}\bigl(\mathsf{blind}_i,\ x(C_{1,i}),\
x(C_{2,i}),\ y(C_{1,i}),\ y(C_{2,i})\bigr),
\]
\begin{sloppypar}
where $d$ is the vote decision (\texttt{src/gadgets/vote\_commitment.rs},
lines 42--55; \texttt{src/shares\_hash.rs}, lines 39--61). The
$y$-coordinates are included deliberately: an $x$-coordinate identifies a
point only up to sign, and hashing both coordinates blocks ciphertext
sign-malleability. The leading tag $1$ separates vote commitments from VANs
(tag $0$) in the shared tree and is baked into the verification key, so an
honest verifier cannot be driven to confuse the two leaf kinds
(\texttt{src/gadgets/vote\_commitment.rs}, lines 11--13). A cast message
reveals $\mathsf{vc}$, the VAN nullifier, the successor VAN, and
$\mathsf{pid}$ --- neither the decision nor the weight. Note that the cast
publicly links $\mathsf{vc}$ to $\nfsym^{\mathsf{van}}$, so per-vote weight
privacy rests on the encryption, not on unlinkability.
\end{sloppypar}

\begin{remark}[Quantization gap, propagated]\label{rem:vote-quant}
The sum condition ties $\sum_i v_i$ to the ballot count $n_b$ minted at
delegation, so its soundness inherits the delegation circuit's
quantization relation: the constraint
$n_b \cdot 12{,}500{,}000 + \mathsf{rem} = v_{\mathrm{total}}$ with
$\mathsf{rem} < 2^{24}$ is strictly weaker than floor division and admits a
one-ballot \emph{under}-claim for roughly $34\%$ of note values; an
over-claim is impossible (Proposition~\ref{thm:vote-underclaim} in
\S\ref{sec:vote-weight}; \texttt{src/params.rs}, lines 9--13;
\texttt{src/delegation/README.md}, \S 8). This is the protocol's one open
problem --- the public design prose states exact floor division, the code
does not enforce it.
\end{remark}

\subsection{The Election Authority: a per-round Joint--Feldman DKG}
\label{sec:vote-dkg}

The tally requires a key that opens aggregates; any party holding that key
alone could equally open each accumulated share the moment it lands on
chain. The protocol therefore splits the round key across the bonded
validator set by a Joint--Feldman distributed key generation, and re-runs
the ceremony every round, so compromise of one round's key reaches no other
round and the validator set may change between rounds.

Each ceremony validator --- registered through a whitelisted Pallas-key
path --- samples a secret polynomial of degree $t - 1$ over the Pallas
scalar field, posts its $t$ Feldman coefficient commitments (Pallas points)
on chain, and delivers to every other validator its evaluation share,
encrypted under ECIES: an ephemeral Pallas Diffie--Hellman, the symmetric
key $\mathrm{SHA\text{-}256}(E \,\|\, S_x)$ over the compressed ephemeral
point and the shared-secret $x$-coordinate, and ChaCha20-Poly1305 with a
zero nonce --- safe because each envelope uses a fresh ephemeral key
(\texttt{crypto/ecies/ecies.go}, lines 21--33). The chain validates point
well-formedness and the commitment count
(\texttt{x/vote/keeper/msg\_server\_ceremony.go}, lines 137--149). When the
$n$-th contribution arrives, the chain sums the commitment vectors
componentwise; the round key $\mathsf{ea_{pk}}$ is the combined constant
term, and each validator's verification key $\mathsf{VK}_i$ is the combined
commitment polynomial evaluated at $i$ in the exponent
(\texttt{x/vote/keeper/msg\_server\_ceremony.go}, lines 224--269). No dealer
ever holds the round secret: it exists only as the sum of shares.

The parameters are fixed in code: reconstruction threshold
$t = \max(2, \lfloor (2n+2)/3 \rfloor)$ for $n \geq 2$ validators (the
degenerate $n = 1$, $t = 1$ case exists for local testing), ceremony
finalization quorum $\lceil 4n/5 \rceil$ acknowledgements --- stricter than
$t$ for most $n$ --- with non-ackers stripped from the ceremony and
non-participants jailed (\texttt{x/vote/keeper/keeper\_ceremony.go}, lines
46--79; \texttt{keeper\_ceremony\_jailing.go}).

The known key-bias attack on plain Joint--Feldman applies: contributions
arrive sequentially through the block-proposer path, so the last
contributor sees every prior commitment and can steer the combined key to
a chosen point. The repository's design document concedes the bias and
argues it is harmless here --- the shifted key reveals nothing about the
round secret under the discrete-log assumption, and the key serves only
ElGamal encryption, whose IND-CPA security holds for any valid key ---
citing Gennaro et al.\ (2007) for the security of Joint--Feldman threshold
decryption despite a biased key (\texttt{vote-sdk/docs/tss-ceremony.md},
``Why single-phase''). The argument is itself designed but unspecified: a
repository document by the designers, not a normative artifact or
independent analysis.

\begin{remark}[Threshold collusion]\label{rem:vote-ea-collusion}
Any $t$ colluding validators can reconstruct the round secret offline and
decrypt every individual on-chain ciphertext; summing one vote's sixteen
shares then yields that vote's exact per-decision weight, which the cast
message already links in public to a VAN nullifier. The linkage stops at the
VAN layer --- the governance nullifier's keyed PRF keeps Zcash notes and
addresses out of reach. No published source claims per-vote weight privacy
against a threshold coalition; a specification would have to state this
limit explicitly (inference from \texttt{crypto/elgamal/bsgs.go} and
\texttt{proto/svote/v1/tx.proto}, lines 88--99; recorded 2026-07-08).
\end{remark}

\subsection{Tally: chain-verified threshold decryption}
\label{sec:vote-decrypt}

\paragraph{Accumulation.} On each accepted reveal, the chain adds the
$64$-byte ciphertext componentwise into the accumulator keyed by the triple
(round, proposal, decision), validating well-formedness and rejecting
identity components (\texttt{x/vote/keeper/keeper\_tally.go}, lines 33--74).
By linearity, the accumulator encrypts the summed ballots of all revealed
shares for that cell. Nothing decrypts until the round enters its tallying
phase.

\paragraph{Partial decryptions.} During tallying, each ceremony validator's
node injects --- through the block-proposer path; the mempool route is
rejected --- one submission covering every non-empty accumulator: the
partial decryption $D_i := [s_i]\,C_1$ under its DKG share $s_i$, with a
Chaum--Pedersen DLEQ proof that
$\log_G \mathsf{VK}_i = \log_{C_1} D_i$. The proof is the pair $(e, z)$; the
verifier recomputes $R_1 = [z]\,G - [e]\,\mathsf{VK}_i$ and
$R_2 = [z]\,C_1 - [e]\,D_i$ and accepts iff
\[
e = H\bigl(\mathtt{svote\text{-}pd\text{-}dleq\text{-}v1} \,\|\, G \,\|\,
\mathsf{VK}_i \,\|\, C_1 \,\|\, D_i \,\|\, R_1 \,\|\, R_2\bigr),
\]
with $H$ a BLAKE2b-256 hash-to-scalar; the tag separates this proof family
from the tally-level DLEQ family (\texttt{svote-dleq-v1}), barring
cross-family replay (\texttt{crypto/elgamal/dleq.go}, lines 30--145). The
verification key $\mathsf{VK}_i$ is re-derived on the fly from the round's
combined Feldman commitments rather than stored
(\texttt{x/vote/keeper/msg\_server\_tally\_decrypt.go}, lines 257--269).
The DLEQ is what makes the tally trustless against its own operators: one
forged $D_i$ would silently shift every total. Duplicate submissions per
validator are rejected, and each submission must cover every non-empty
accumulator, so a proposer cannot inject a partial set that later wedges the
tally (\texttt{msg\_server\_tally\_decrypt.go}, lines 199--294).

\paragraph{Final verification.} The block proposer alone may submit the
claimed plaintext totals. For each entry the chain Lagrange-combines at
least $t$ stored partials at zero,
\[
\sum_i \lambda_i\, D_i = \Bigl[\sum_i \lambda_i\, s_i\Bigr] C_1 =
[s]\,C_1,
\]
recovering the round secret only in the exponent
(\texttt{crypto/shamir/partial\_decrypt.go}, lines 37--71), and checks
$C_2 - [s]\,C_1 = [v]\,G$ against the claimed total $v$; an empty
accumulator forces $v = 0$; any $v \geq 2^{28}$ is rejected outright (the
constant \texttt{TallyBSGSBound}, an exclusive bound;
\texttt{x/vote/types/keys.go}, lines 33--35). A completeness check requires
the entries to cover every non-empty accumulator --- without it a malicious
proposer could finalize with results omitted
(\texttt{msg\_server\_tally\_decrypt.go}, lines 102--177). Only then does
the round finalize.

\paragraph{Recovering the logarithm.} The chain verifies by re-encoding
$[v]\,G$, so only the proposer must solve discrete logarithms, by
baby-step/giant-step under $N = 2^{28}$: a table of $m = \lceil \sqrt{N}
\rceil = 2^{14} = 16{,}384$ baby steps $j \mapsto [j]\,G$, giant steps
subtracting $[m]\,G$, a hit at $(i, j)$ giving $v = i m + j$ within at most
$m + 1$ iterations (\texttt{crypto/elgamal/bsgs.go}, lines 11--110). The
bound is ample: the full $21$-million-ZEC supply quantizes to
$1.68 \times 10^{8}$ ballots at $12{,}500{,}000$ zatoshi per ballot
(\texttt{src/params.rs}, line 14), below $2^{28} \approx 2.68 \times
10^{8}$, so no honest aggregate is ever truncated.

\paragraph{Two flags.}
\begin{sloppypar}
Totals are denominated in \emph{ballots}; the
protobuf comment on the tally field says zatoshi and is wrong --- the wallet
multiplies by $12{,}500{,}000$ for display
(\texttt{proto/svote/v1/tx.proto}, lines 111--116, against
\texttt{zashi-android} \texttt{TallyResults.kt}, lines 44--56, commit
\texttt{b4edd6f6}). And the tallying phase carries a timeout of $21{,}600$
seconds (six hours), after which the round finalizes \emph{without}
results (\texttt{x/vote/types/keys.go}, lines 28--31) --- an availability
caveat, not an integrity one, but a specification must decide whether a
resultless finalization is acceptable.
\end{sloppypar}

\paragraph{What a specification must pin down.} Until zips PR~\#1200
merges, the commit-pinned sources above are the only citable definition of:
the Poseidon preimage layouts and the domain-tag registry (currently pinned
by unit tests); the three nullifier formulas, their set scoping, and their
key encodings; the decrement lookup table and the sentinel-bit convention;
the share count, the per-share range bound, and their interaction with
partial reveals; the DKG parameter formulas, ceremony state machine, and
jailing policy; the DLEQ transcript formats and tags; and the BSGS bound and
the unit denomination of totals as consensus parameters.

\section{Trust model and verdict}\label{sec:vote-trust}

The preceding sections constructed the protocol's mechanisms from
commit-pinned source (Table~\ref{tab:vote-sources}); this section assembles
what those mechanisms do \emph{not} guarantee. We first map every party the
protocol trusts and state exactly what each is trusted with, then classify
the volume's claims into the three bins of
\S\ref{sec:vote-source-situation}, and close with the verdict: what a
formal specification would have to pin down before the protocol's
guarantees become checkable, and what the protocol delivers today against
the instruments that ran the same June 2026 question. File citations below
are at the commits of Table~\ref{tab:vote-sources}; wallet files are at
\texttt{zashi-android} @ \texttt{b4edd6f6}.

Every trust relation in this map is itself designed but unspecified: each
is concrete, open-source behaviour at the pinned commits, and none is
stated in a normative artifact. The inventory of
\S\ref{sec:vote-inventory} records the full classification.

\subsection{The trust map}\label{sec:vote-trust-map}

\paragraph{Helper servers hold the vote decision and the linkage.} The
payload a wallet sends to each chosen helper --- struct
\texttt{SharePayload} in \texttt{zcash\_voting/src/types.rs} --- carries
the plaintext \texttt{vote\_decision}, the vote commitment's
\texttt{tree\_position} in the per-round tree, the \texttt{shares\_hash},
all sixteen ciphertexts, the sixteen share commitments, and the revealed
share's blinding factor (\texttt{primary\_blind}); the helper builds the
Merkle path and the share-reveal proof server-side and submits
\texttt{MsgRevealShare} itself (\texttt{vote-sdk/internal/helper/%
processor.go}). A helper therefore learns, before any reveal, how the
voter voted, and holds precisely the share-to-commitment linkage that the
on-chain data blinds; composed with the public pairing of
\texttt{vote\_commitment} and \texttt{van\_nullifier} in
\texttt{MsgCastVote} (\texttt{proto/svote/v1/tx.proto}), that linkage ties
the decision to the vote's delegated authority instance. This is a broader
grant than a timing-relay framing suggests. What a helper never receives
is weight: the plaintext share values and the ElGamal randomness are marked
\texttt{SECRET} in the client library and stripped from the wire type
(\texttt{types.rs}, \texttt{EncryptedShare} versus
\texttt{WireEncryptedShare}). The timing decorrelation itself is client
policy, not a protocol guarantee: the wallet posts each share to half of
the configured helpers rounded up, each with an independently random
future \texttt{submit\_at} scheduled no later than a last-moment buffer
--- two fifths of the round duration, capped at six hours --- before the
vote deadline (\texttt{zcash\_voting/src/share\_policy.rs}); a censoring
helper causes
partial weight loss, mitigated only by the multi-server fan-out and by a
single-share fallback described below.

\paragraph{Any $t$ colluding validators decrypt every posted ciphertext.}
Function \texttt{ThresholdForN} sets the reconstruction threshold
$t = \max(2, \lfloor (2n+2)/3 \rfloor)$ for $n$ ceremony validators
(\texttt{vote-sdk/x/vote/keeper/keeper\_ceremony.go}). Any $t$ of them can
reconstruct the round's ElGamal secret key offline from their dealt shares
and decrypt every \texttt{enc\_share} on the chain --- not merely the
aggregates the tally publishes. The decision attached to each ciphertext
is public at reveal time in any case (\texttt{MsgRevealShare} carries
\texttt{vote\_decision}); what threshold collusion adds is the share
values. In the single-share last-moment mode the exposure is immediate:
the client builds a payload for share~0 only, ``which carries all the
voting weight'', and the remaining fifteen zero-value shares are never
revealed (\texttt{zcash\_voting/src/vote\_commitment.rs}), so one
decryption yields one vote's exact weight. In the sixteen-share mode,
recovering a per-vote weight additionally requires grouping shares by
vote --- a linkage the chain withholds (share nullifiers are blinded), the
helpers of the previous paragraph possess outright, and the client-side
timing policy merely obfuscates. The linkage stops at the Vote Authority
Note layer in every case: the governance nullifier is keyed by the
holder's nullifier-deriving key \nksym{}, so mainnet notes and addresses
stay unlinkable even after a later spend
(\texttt{voting-circuits/src/delegation/circuit.rs}, condition~14). We
flag the analytical gap explicitly: neither the public documentation nor
the source claims --- or denies --- per-vote weight privacy against a
colluding Election Authority threshold. The limit stated here is our
adversarial inference from \texttt{crypto/elgamal/bsgs.go}, the message
surface, and the share-range condition of the vote-proof circuit; the
property is absent from every published analysis.

\paragraph{The chain is closed and its lifecycle is coordinator-gated.}
The vote chain enforces a positive-security message whitelist at the ante
handler: only enumerated message types pass, standard bank transfers are
disabled --- explicitly ``because unrestricted transfers would let anyone
accumulate enough stake to create a validator, bypassing the controlled
validator set'' (\texttt{vote-sdk/proto/svote/v1/tx.proto}, comment on
\texttt{MsgAuthorizedSend}) --- and post-genesis validator creation is
disallowed outside the registered-Pallas-key path
(\texttt{vote-sdk/app/ante\_whitelist.go}). Round lifecycle, privileged
funding, and software upgrades are threshold-gated \emph{coordinator
actions} proposed and approved by a vote-manager set
(\texttt{MsgProposeCoordinatorAction}, \texttt{MsgApproveCoordinatorAction});
the default genesis ships a single vote manager, and the coordinator
threshold defaults to one (\texttt{vote-sdk/x/vote/module.go},
\texttt{DefaultVoteManagerAddresses}; \texttt{tx.proto},
\texttt{MsgUpdateVoteManagers}).
%% TODO: production vote-manager set, coordinator threshold, and validator
%% identities not establishable from the pinned sources; needs the deployed
%% chain's genesis or a valargroup announcement.
The same coordinators post the snapshot: the eligibility anchors
\texttt{nc\_root} and \texttt{nullifier\_imt\_root} arrive in
\texttt{MsgCreateVotingSession} as client-fetched values --- the
note-commitment root recomputed from a lightwalletd frontier, the
nullifier root taken from the operator's PIR service
(\texttt{vote-sdk/api/snapshot.go}) --- and the chain stores them without
verifying that they describe any real Zcash state; it runs no light
client (\texttt{x/vote/keeper/msg\_server.go}). Honest wallets' proofs
fail against fabricated roots, and anyone can audit the posted roots
after the fact, but nothing in the protocol prevents a coordinator quorum
from pinning a snapshot of its own invention.

\paragraph{The wallet vendor's configuration is the authentication root.}
A voter's assurance that they are voting in a genuine round reduces to a
static configuration file bundled with the wallet: fetched from a
commit-pinned URL (\texttt{valargroup/token-holder-voting-config} @
\texttt{2785311d}) and verified against a hard-coded SHA-256 digest, it
names the ed25519 keys that must sign the dynamic per-round
configuration --- helper endpoints, PIR endpoints, and each round's
Election Authority key (\texttt{StaticVotingConfig.kt},
\texttt{BUNDLED\_PINNED\_SOURCE}). The round authenticator rejects a round
whose on-chain \texttt{ea\_pk} disagrees with the signed entry
(\texttt{RoundAuthenticator.kt}, status \texttt{EA\_PK\_MISMATCH}) --- the
wallet's only defence against an attacker-substituted encryption key,
under which every cast share would be attacker-decryptable. An on-chain
endorsement registry (\texttt{MsgEndorseRound}) is secondary. The trust
root of the whole client-side protocol is therefore the custodians of the
config-signing keys together with the application distribution channel;
neither is named in any published artifact.

\paragraph{Tally liveness admits finalization without results.} The tally
verification itself is real: validators' nodes auto-inject partial
decryptions whose Chaum--Pedersen DLEQ proofs are checked on-chain against
Feldman-derived verification keys, and the proposer-only
\texttt{MsgSubmitTally} is accepted only if every claimed total matches
the Lagrange combination of stored partials, covers every non-empty
accumulator, and lies below the discrete-log recovery bound $2^{28}$
(\texttt{vote-sdk/x/vote/keeper/msg\_server\_tally\_decrypt.go};
\texttt{crypto/shamir/feldman.go}; the \texttt{TallyEntry} field
\texttt{decryption\_proof} is dead --- ``reserved for future use'' --- the
binding check is the recomputation). But the threshold $t$ is a liveness
parameter as well as a secrecy parameter: if fewer than $t$ validators
decrypt, the round sits in \texttt{TALLYING} until the timeout ---
$21{,}600$ seconds by default (\texttt{x/vote/types/keys.go},
\texttt{DefaultTallyTimeout}) --- and the end-blocker then finalizes it
with \texttt{tally\_timed\_out = true} and \emph{empty results}
(\texttt{x/vote/module.go}). A finalized round in which every cast vote
remains permanently undisclosed is a protocol outcome, not an excluded
failure mode; the design accepts it as the price of not blocking the
chain on an unreachable threshold.

\paragraph{The tally wire format mislabels its units.} The comment on
\texttt{TallyEntry.total\_value} reads ``Decrypted aggregate (zatoshi)''
(\texttt{vote-sdk/proto/svote/v1/tx.proto:114}). It is wrong: the
accumulated plaintexts are \emph{ballot counts} --- the vote-proof circuit
constrains each vote's shares to sum to the Vote Authority Note's ballot
count, and the wallet converts totals to zatoshi by multiplying by the
ballot divisor $12{,}500{,}000$ (\texttt{TallyResults.kt},
\texttt{toTallyResults}; constant \texttt{BALLOT\_DIVISOR\_ZATOSHI} in
\texttt{VotingCryptoClient.kt}). The recovery bound corroborates the
correct reading: $2^{28}$ ballots is roughly $33.6$ million ZEC ---
comfortably above the monetary supply --- whereas $2^{28}$ zatoshi is
under $2.7$ ZEC. A consumer that trusts the comment under-reports every
result by a factor of $12.5$ million. In a protocol whose only
specification is its implementation, a wrong comment in the wire
definition is not a cosmetic defect; it is an error in the closest thing
the protocol has to a spec.

\subsection{Claim inventory: the three bins}\label{sec:vote-inventory}

\paragraph{Specified.} Only the inherited primitives: \NoteCommit{},
nullifier derivation, \CommitIvk{}, the \Sinsemilla{} Merkle instance,
RedPallas, the ZIP-244 sighash, and ZIP-32 derivation, each specified in
the protocol specification and the ZIPs, constructed in the \emph{Orchard
Guide} (\S``Representation of a note: notes and note commitments'',
\S``Prevention of double-spending: nullifiers'', \S``Keys: capabilities
for spending and for viewing'', \S``Proof of ownership of \emph{some}
valid note: the commitment tree'', \S``Authorisation of the spend:
RedPallas''), and reused unchanged through the
\texttt{unstable-voting-circuits} feature (zcash/orchard @
\texttt{c0b6b917}). No governance-specific object reaches this bin.

\paragraph{Designed but unspecified.} Everything this volume constructed:
the five-message ballot surface; the delegation relation with its
note-bundle eligibility and indexed-Merkle-tree non-membership argument;
the coordinator-posted, chain-unverified snapshot roots; the Vote
Authority Note and its synthetic-action authorization --- including the
asymmetry that the chain verifies the delegation signature over a
\emph{client-provided} sighash, relying on the circuit's $\rho$-chain for
binding, while the cast-vote sighash is recomputed on-chain
(\texttt{x/vote/ante/validate.go}); all four double-vote layers; the
sixteen-share lifted-ElGamal encoding, vote commitment, and bitmask
decrement; the helper protocol and the trust it grants; the per-round
Joint--Feldman DKG with threshold $t = \max(2, \lfloor (2n+2)/3 \rfloor)$
and finalization quorum $\lceil 4n/5 \rceil$; the chain-verified tally,
its $2^{28}$ recovery bound, and the timeout path; the PIR eligibility
subsystem; and every trust relation of \S\ref{sec:vote-trust-map}. Each is
concrete and open-source at the pinned commits; none is normative. The
would-be specifications --- zcash/zips PRs \#1200 and \#1198, and the
prior-art PR \#853 --- were all open and unmerged as of 2026-07-08.

\paragraph{Open problem.} One: the ballot-quantization relation. The
delegation circuit constrains
$\mathit{num\_ballots} \cdot 12{,}500{,}000 + \mathit{remainder} =
v_{\mathrm{total}}$ with $\mathit{remainder} < 2^{24}$, which is strictly
weaker than the floor division the public documentation states: the pair
$(q-1,\, r+12{,}500{,}000)$ also satisfies every constraint whenever the
honest remainder obeys
$r < 2^{24} - 12{,}500{,}000 = 4{,}277{,}216$ --- about $34.2\%$ of
residues --- and the honest ballot count is at least two, while the
opposite deviation fails the range check, so a
prover can under-claim exactly one ballot and can never over-claim
(\texttt{voting-circuits/src/delegation/circuit.rs};
\texttt{src/delegation/README.md}~\S8, which states the window, proves
over-claim impossible, and lists unimplemented tightening constraints).
The deviation is self-harming, but it is the volume's one confirmed
disagreement between the code and the code's own public description, and
the exact relation a specification must choose --- floor division or the
shipped weaker relation --- is unresolved.

\subsection{Verdict}\label{sec:vote-verdict}

\paragraph{What a specification must pin down.} For the protocol's
guarantees to become checkable --- by an auditor, by a second
implementation, or by this volume --- a normative artifact would need to
fix, at minimum:
\begin{enumerate}[itemsep=1pt]
\item the three circuit relations in full: public-input layouts, every
  condition, and the exact quantization rule --- either closing the
  under-claim window or normatively accepting the shipped relation;
\item every governance-object formula: the four nullifiers and the
  nullifier-domain derivation, the Vote Authority Note commitment
  (with the argument that its $x$-only encoding is bound elsewhere), the
  vote commitment, and the shares hash;
\item the snapshot: how \texttt{nc\_root} and the nullifier root are
  derived, and \emph{who verifies them} --- at present, on-chain, nobody;
\item the authorization model: which sighashes the chain recomputes,
  which it accepts from the client, and why the delegation path is safe;
\item the Election Authority: the threshold and quorum formulas, key
  registration and rotation, and --- the largest analytical gap --- an
  explicit confidentiality claim stating against whom per-vote weight
  privacy holds;
\item the helper interface: the exact payload, the trust granted by it,
  and the status of the timing policy as guarantee or convention;
\item tally semantics: the verification equation, the unit of
  \texttt{total\_value} (ballots), and the meaning of a timed-out round;
\item independent verification: a second implementation of at least the
  verifiers, breaking the single-crate tautology of
  \S\ref{sec:vote-source-situation}.
\end{enumerate}

\paragraph{What the protocol delivers today.} Measured against the
instruments beside it (\S\ref{sec:vote-lineage}), the deliverable is
real but conditional. Against ZCAP, the protocol opens participation to
every Orchard coinholder in proportion to holdings, with no identified
panel and no registrar, and it hides identity and weight from public
observers --- where ZCAP offers the converse trade: an identified panel
on mature, independently analysed voting software. Against the prior
coin-weighted instrument, hhanh00's \texttt{zcash-vote}, it adds
threshold decryption in place of a single tallying authority, blinded
delegation through the Vote Authority Note, hardware signing,
PIR-private eligibility queries, and wallet-native deployment --- at the
price of the trust map of \S\ref{sec:vote-trust-map}: a permissioned
chain, a coordinator quorum, helper servers, and a vendor-held
configuration root. Integrity of the aggregate totals is the strongest
property: the tally check is executed on-chain and every input to it is
public, so any observer of the vote chain can replay it, and a
fabricated snapshot is auditable after the fact. Individual privacy is
the weakest: it holds against outside observers, degrades to helper
honesty for the linkage of decisions to votes, and degrades to the
honesty of all but $t-1$ validators for individual weights --- a limit no
published analysis states. On that evidence the protocol is fit for
what it was used for in June 2026: a non-binding sentiment poll with
publicly recomputable totals. Binding governance would additionally
require the specification above, an independent verifier, and a
confidentiality analysis of the Election Authority --- none of which
existed as of 2026-07-08.

\end{document}
